% Generated draft: requires mathematical and source review.
\documentclass[11pt]{article}
\usepackage[margin=25mm]{geometry}
\usepackage{fontspec}
\setmainfont{Latin Modern Roman}
\usepackage{amsmath,amssymb,mathtools,unicode-math}
\setmathfont{Latin Modern Math}
\usepackage{enumitem,needspace,xurl,hyperref}
\setlength{\parindent}{0pt}
\setlength{\parskip}{5pt}
\setlength{\emergencystretch}{3em}
\newcommand{\N}{\mathbb{N}}
\newcommand{\Z}{\mathbb{Z}}
\newcommand{\Q}{\mathbb{Q}}
\newcommand{\R}{\mathbb{R}}
\newcommand{\C}{\mathbb{C}}
\newcommand{\F}{\mathbb{F}}
\newcommand{\A}{\mathbb{A}}
\newcommand{\PP}{\mathbb{P}}
\newcommand{\E}{\mathbb{E}}
\newcommand{\eps}{\varepsilon}
\newcommand{\dd}{\,\mathrm{d}}
\begin{document}
\section*{UnsolvedMath ID: 800006}
\textbf{Source problem code:} catalogue entry
\subsection{Definitions and mathematical statement}
A semigroup is a nonempty set with an associative binary operation. A semigroup variety is a class of semigroups defined by identities. Let \(\mathcal L\) be the lattice of all semigroup varieties, ordered by inclusion. For \(\mathcal V\in\mathcal L\), define its set of covers by
\[
\operatorname{Cov}_{\mathcal L}(\mathcal V)=\{\mathcal W\in\mathcal L:\mathcal V\subsetneq\mathcal W\text{ and there is no }\mathcal U\in\mathcal L\text{ with }\mathcal V\subsetneq\mathcal U\subsetneq\mathcal W\}.
\]
For a source-intended index \(n\) for which the displayed power is meaningful as a semigroup word, put
\[
\mathcal E_n=\{S:S\text{ is a semigroup and }a^2=a^{n+2}\text{ for every }a\in S\},\qquad
\mathcal L_n=\{\mathcal V\in\mathcal L:\mathcal V\subseteq\mathcal E_n\}.
\]
Covers in \(\mathcal L_n\) are defined by the same strict-inclusion condition, with all intervening varieties required to belong to \(\mathcal L_n\).

\textbf{Problem (S. Kublanovsky).} Describe
\[
\operatorname{Cov}_{\mathcal L}(RS_n)
\quad\text{and}\quad
\operatorname{Cov}_{\mathcal L_n}(RS_n).
\]
The second expression presupposes \(RS_n\subseteq\mathcal E_n\). The supplied record does not define the semigroup variety \(RS_n\), specify the range of \(n\), or state hypotheses establishing this containment. Thus this is a conditional formulation of the requested classification; the target varieties cannot be identified precisely from the supplied record alone.

The source reports that results labeled [Kubl5] imply that \(\operatorname{Cov}_{\mathcal L_n}(RS_n)\) is finite. It also reports, citing [Kubl4], that every subvariety of \(RS_1\) has finitely many covers. The supplied excerpt does not explicitly specify the ambient lattice for this latter assertion. These are retained as source-reported background statements, without independent verification.
\subsection{Short English statement}
Describe all varieties immediately above \(RS_n\), both among all semigroup varieties and among those satisfying \(x^2=x^{n+2}\). The supplied record leaves \(RS_n\) and the range of \(n\) undefined, so the classification target requires clarification. Its finiteness statements are supplied background claims.
\subsection{Sources}
\begin{itemize}
\item Local UnsolvedMath registry: permanent ID \(800006\), source problem code \texttt{catalogue entry}, title \emph{The Lyapin's Notebook --- Problem 6}. The registry labels the problem open; this is a supplied status, not an independently verified assessment.
\item Bershadsky, Kublanovsky, Mashevitzky, \emph{The Lyapin's Notebook: Unsolved Problems in Semigroup Theory} (2026), Problem 6, attributed in the statement to S. Kublanovsky. Supplied source URL: \url{https://arxiv.org/abs/2604.04763}. The registry records access on 2026-10-16 and extraction by \texttt{arxiv\_latex}; these are supplied provenance details, not a claim of access or verification here. The registry's proposer field lists Bershadsky, Kublanovsky, Mashevitzky and its proposed-year field gives 2026.
\item [Kubl5]: cited in the supplied statement for finiteness of the covers in the lattice satisfying \(x^2=x^{n+2}\). No bibliographic entry or separate URL is supplied.
\item [Kubl4]: cited in the supplied statement for finiteness of the covers of every subvariety of \(RS_1\). No bibliographic entry or separate URL is supplied.
\end{itemize}
\end{document}
